Across \NCheckpoints{} checkpoints and \NExecEpisodes{} episodes, \NUninformative{} of the template
pairs turned out to render \emph{this} conversation identically despite differing on the probe set,
because the edit only touches a branch this conversation does not take. Those pairs vary nothing and
are excluded: an identical input producing an identical output is not evidence that the edit was
harmless. That leaves \NFamily{} informative comparisons, and both outcome families are
Holm-corrected across all \NFamily{}.

\paragraph{The floor first, because it disqualifies the obvious metric.} Running the same template
twice, byte-identical prompt and settings, changes the reply on \MinNoiseFloor\% of tasks for
\ShortMistral{} and \ShortPhi{} but on \MaxNoiseFloor\% for \ShortDS{}. Response divergence is
therefore useless as a drift signal on its own, since on \ShortDS{} it is near-saturated before any
template changes at all. What survives replication is the executed outcome, and that is what we test.

\paragraph{A template edit can make the request impossible.} \ShortMistral's original template
raises on a system message. All \MistralRenderErr{} tasks fail before the model is reached, giving
resolution \ResMistralA\% and contract compliance \ConMistralA\%. A later edit to the same
repository, on the same weights, accepts system messages: resolution \ResMistralB\% and contract
\ConMistralB\%. This is the one resolution comparison that survives Holm across all \NFamily{} tests
($p_{\text{Holm}}=\MistralPHolmRes$; the contract swing also survives, at
$p_{\text{Holm}}\,{\MistralPHolmCon}$), and it is categorical rather than statistical. A pipeline
pinned to the earlier revision cannot send a system prompt at all, and one written against the later
revision fails completely if it pulls the earlier one.

\paragraph{A template edit can halve what a prompt suite asserts, and a later edit can restore it.}
This is the strongest execution-grounded effect we observe, and it is a one-repository case study:
it establishes that the failure occurs, not how often. On \ShortPhi{}, contract compliance runs
\ConPhiA\%, then \ConPhiB\%, then \ConPhiC\% across the
repository's three templates, while resolution runs \ResPhiA\%, \ResPhiB\%, \ResPhiC\%. Both
contract transitions survive correction, down at $p_{\text{Holm}}\,{\PhiPHolmConDown}$ and back up
at $p_{\text{Holm}}\,{\PhiPHolmConUp}$. The size of the swing is the operational point rather than the
p-value: compliance halves at one edit and is restored at the next, a movement far outside the
checkpoint's own replication floor. The middle template is the commit whose own message describes
replacing the system role with the user role, which is what a halved compliance rate on a
system-carried output contract looks like from the outside. The resolution swing over the same two
transitions is large but does not survive correction across \NFamily{} tests
($p_{\text{Holm}}=\PhiPHolmResDown$ and \PhiPHolmResUp); with only \PhiDiscDown{} and \PhiDiscUp{}
discordant tasks the smallest p-value these exact tests can produce cannot clear a family of
\NFamily{}, so the design is underpowered for these transitions and we report them as suggestive
rather than as established or refuted.

\paragraph{This is where a then-versus-now check fails.} Comparing \ShortPhi's first template with
its last gives $p_{\text{Holm}}=\PhiFirstLastRes$ on resolution and a contract rate back to roughly
where it started. The repository broke and then repaired itself. A check that compares only the
current template against the one it was pinned at sees nothing, while what a consumer lived through
was the window in between.

\paragraph{Drift is often harmless.} \ShortDS{} ships \NVersDSCoder{} distinct
templates and none of the \NDSPairs{} informative comparisons between them moves resolution detectably, the
smallest being $p_{\text{Holm}}=\DSMinP$, with resolution staying between \ResDSCoderA\% and
\ResDSCoderF\%. Of the \NFamily{} informative comparisons, \NSigRes{} survives correction on
resolution and \NSigContract{} on contract compliance. The claim is not that template edits usually
break things. It is that some do, that the ones that do cannot be told from the ones that do not
without running the model, and that nothing in the tooling available when this study was run tells
a consumer which kind has just landed.

\paragraph{Two robustness checks, because both cut against us if they fail.} First, the excluded
uninformative pairs are put back and everything is re-corrected over all \NAllPairs{} pairs. This is
the \emph{more} conservative direction, since each restored pair contributes $p=1$ by construction
and enlarges the family, and the counts are unchanged: \NAllSigRes{} on resolution and
\NAllSigCon{} on contract. Second, each effect is read against its own checkpoint's replication
discordance. \ShortPhi{} and \ShortMistral{} show \FloorDiscPhi{} and \FloorDiscMistral{}
discordant tasks when a
template is replayed against itself, against \PhiDiscDown{} and \PhiDiscUp{} for the two \ShortPhi{}
transitions and \MistralDisc{} for \ShortMistral{}, so those effects sit well outside their own
noise. \ShortMistral's floor is trivially clean, since its replicated arm fails to render every
task; its effect rests on the categorical contrast above, not on that floor. \ShortQwen{} is the cautionary case in the other direction: replaying an identical prompt
moves its resolution enough that the replication arm would itself test significant
($p=\QwenFloorP$) against a zero-discordance null. Any single comparison on a checkpoint like that
should be disbelieved, which is precisely why the floor is measured rather than assumed.
